% Part: computability
% Chapter: tm-computations
% Section: church-turing-thesis

\documentclass[../../../include/open-logic-section]{subfiles}

\begin{document}

\olfileid{tur}{mac}{ctt}
\olsection{ચર્ચ--ટ્યુરિંગ માન્યતા}

ટ્યુરિંગ મશીનો અસરકારક પ્રક્રિયાના ખ્યાલનું ચોક્કસ
ગણિતીય સ્વરૂપ આપવાનાં છે. ટ્યુરિંગના મતે, જે કોઈ
અસરકારક પ્રક્રિયા અને ટ્યુરિંગ મશીન બંનેના ખ્યાલને
સમજે તેને સહજ રીતે લાગશે કે અસરકારક પ્રક્રિયા વડે
જે કંઈ કરી શકાય તે ટ્યુરિંગ મશીન વડે પણ કરી શકાય.
આ દાવાને એ હકીકત ટેકો આપે છે કે અસરકારક પ્રક્રિયાના
ખ્યાલનાં અન્ય તમામ પ્રસ્તાવિત ચોક્કસ સ્વરૂપો ટ્યુરિંગ
મશીનના ખ્યાલ સાથે વિસ્તારની દૃષ્ટિએ સમકક્ષ નીવડે
છે---એટલે તેઓ બરાબર એ જ વિધેયોનો ગણ ગણી શકે છે.
આ દાવો \emph{ચર્ચ--ટ્યુરિંગ માન્યતા} કહેવાય છે.

\begin{defn}[ચર્ચ--ટ્યુરિંગ માન્યતા]
\emph{ચર્ચ--ટ્યુરિંગ માન્યતા} કહે છે કે અસરકારક
પ્રક્રિયા વડે સંગણનીય હોય તે બધું ટ્યુરિંગ-સંગણનીય છે.
\end{defn}

ચર્ચ--ટ્યુરિંગ માન્યતાનો બે રીતે આધાર લેવાય છે. પ્રથમ
રીત આળસને બહાનું આપે છે. ધારો કે આપણને કંઈક ગણવાની
અસરકારક પ્રક્રિયાનું વર્ણન, માનો ``આભાસી કોડ''માં,
મળ્યું છે. તો આપણે ખરેખર ટ્યુરિંગ મશીન બનાવ્યું ન
હોય તોપણ એ જ વિધેય કોઈ ટ્યુરિંગ મશીન ગણે છે એવો
દાવો યોગ્ય ઠેરવવા ચર્ચ--ટ્યુરિંગ માન્યતા લઈ શકીએ.

ચર્ચ--ટ્યુરિંગ માન્યતાનો બીજો ઉપયોગ તત્ત્વચિંતનની
દૃષ્ટિએ વધુ રસપ્રદ છે. કેટલાંક વિધેયો ટ્યુરિંગ મશીન
વડે ગણી શકાતા નથી તે બતાવી શકાય છે. તેના પરથી,
ચર્ચ--ટ્યુરિંગ માન્યતા વાપરીને, કોઈ પણ પ્રકારની પ્રક્રિયા
વડે તેમને અસરકારક રીતે ગણી શકાતા નથી એવો નિષ્કર્ષ
કાઢી શકાય છે. કારણ કે એવી પ્રક્રિયા હોય, તો
ચર્ચ--ટ્યુરિંગ માન્યતા મુજબ તેને માટે ટ્યુરિંગ મશીન
હોવું જોઈએ. તેથી જો તેને ગણતું ટ્યુરિંગ મશીન નથી
એવું સાબિત કરી શકીએ, તો અસરકારક પ્રક્રિયા પણ
હોઈ શકે નહીં. ખાસ કરીને, કહેવાતી અટકવાની સમસ્યા
ટ્યુરિંગ મશીનોથી જ નહીં પરંતુ કોઈ પણ અસરકારક રીતે
ઉકેલી શકાતી નથી એવો દાવો કરવા ચર્ચ--ટ્યુરિંગ
માન્યતા વપરાય છે.

\end{document}
